\documentclass{beamer}
\usetheme{Madrid}

\title{Gig Economy and Changing Labor Market Dynamics}
\subtitle{Contemporary Economic Issues: ECON 204}
\author{Muhebullah Karimzada}
\date{Winter 2025}

\begin{document}

% Title Slide
\begin{frame}
    \titlepage
\end{frame}


\begin{frame}{What is the Gig Economy?}
    \begin{itemize}
        \item The \textbf{gig economy} refers to a labor market characterized by the prevalence of short-term, flexible, and freelance work rather than long-term employment.
        \item Workers in the gig economy typically engage in \textit{temporary} or \textit{part-time jobs}, often facilitated by digital platforms or mobile apps.
        \item Gig workers are generally considered independent contractors, not employees, which means they may not receive traditional benefits like health insurance or retirement plans.
    \end{itemize}
\end{frame}

\begin{frame}{Examples of Gig Economy Work in Canada}
    \begin{itemize}
        \item \textbf{Ride-Sharing Services:} \textit{Uber} and \textit{Lyft} allow individuals to work as drivers on a flexible schedule.
        \item \textbf{Food Delivery:} Platforms such as \textit{SkipTheDishes} and \textit{DoorDash} enable people to deliver food as independent contractors.
        \item \textbf{Freelancing:} Websites like \textit{Upwork} and \textit{Freelancer} connect freelancers with clients seeking services in various fields such as writing, graphic design, and programming.
        \item \textbf{Task-Based Work:} Platforms like \textit{TaskRabbit} offer individuals the opportunity to perform various tasks such as assembling furniture or running errands for pay.
    \end{itemize}
\end{frame}



% Problem 1 Slide
\begin{frame}{Problem 1: Wage Disparity in the Gig Economy}
    A gig platform has two types of workers: high-skill and low-skill.
    \begin{itemize}
        \item High-skill workers complete 60\% of tasks, earning an average of \$35/hour.
        \item Low-skill workers complete the remaining 40\%, earning \$20/hour on average.
        \item Workers are paid 80\% of what clients pay, and the platform retains the remaining 20\%.
    \end{itemize}
    \textbf{Questions:}
    \begin{enumerate}
        \item Calculate the platform's total revenue and the wage gap between high-skill and low-skill workers if 5,000 tasks are completed in a week.
        \item Assuming a 10\% increase in tasks next week with the same skill distribution, what is the additional revenue earned by the platform?
    \end{enumerate}
\end{frame}

% Solution 1 Slide
\begin{frame}{Solution 1: Wage Disparity in the Gig Economy}
    \textbf{Given Data:}
    \begin{itemize}
        \item High-skill tasks: 60\% of 5,000 = 3,000 tasks.
        \item Low-skill tasks: 40\% of 5,000 = 2,000 tasks.
        \item High-skill wage = \$35/hour; Low-skill wage = \$20/hour.
        \item Workers earn 80\% of client payments.
    \end{itemize}

\end{frame}



% Solution 1 Slide
\begin{frame}{Solution 1: Wage Disparity in the Gig Economy}

    \textbf{Solution:}
    \begin{align*}
        \text{Total payment by clients} &= \frac{\text{Worker earnings}}{0.8} \\
        \text{High-skill revenue} &= \frac{3,000 \times 35}{0.8} = \$131,250 \\
        \text{Low-skill revenue} &= \frac{2,000 \times 20}{0.8} = \$50,000 \\
        \text{Total revenue} &= \$131,250 + \$50,000 = \$181,250 \\
        \text{Wage gap} &= 35 - 20 = \$15/hour.
    \end{align*}
    \textbf{10\% Increase:}
    \begin{align*}
        \text{New total tasks} &= 5,000 \times 1.1 = 5,500 \\
        \text{Additional revenue} &= 10\% \times \$181,250 = \$18,125.
    \end{align*}
\end{frame}





% Problem 2 Slide
\begin{frame}{Problem 2: Labor Supply Elasticity}
    In a city, the supply of gig workers follows the elasticity formula:
    \[ E_s = \frac{\Delta Q_s / Q_s}{\Delta W / W} \]
    \textbf{Given:}
    \begin{itemize}
        \item Initial wage = \$20/hour; New wage = \$22/hour.
        \item Initial workers = 10,000; New workers = 11,500.
        \item Elasticity \(E_s = 1.5\).
    \end{itemize}
    \textbf{Questions:}
    \begin{enumerate}
        \item Verify if the elasticity aligns with the observed data.
        \item Predict the percentage increase in worker supply if wages increase to \$25/hour.
    \end{enumerate}
\end{frame}

% Solution 2 Slide
\begin{frame}{Solution 2: Labor Supply Elasticity}
    \textbf{Verification:}
    \begin{align*}
        \Delta Q_s &= 11,500 - 10,000 = 1,500, \quad Q_s = 10,000 \\
        \Delta W &= 22 - 20 = 2, \quad W = 20 \\
        E_s &= \frac{(1,500 / 10,000)}{(2 / 20)} = \frac{0.15}{0.1} = 1.5.
    \end{align*}
    \textbf{Prediction:}
    \begin{align*}
        \Delta W / W &= \frac{25 - 20}{20} = 0.25 \\
        \Delta Q_s / Q_s &= E_s \times (\Delta W / W) = 1.5 \times 0.25 = 0.375 \\
        \text{Percentage Increase} &= 0.375 \times 100 = 37.5\%.
    \end{align*}
\end{frame}


% Problem 3 Slide
\begin{frame}{Problem 3: Productivity and Labor Surplus}
    A gig platform provides services in urban areas. The marginal productivity of each worker is given by:
    \[ MP = 50 - 2N \]
    where \(N\) is the number of workers, and marginal productivity is measured in dollars.
    \begin{itemize}
        \item The platform pays workers a fixed wage of \$20/hour.
        \item Currently, 10 workers are employed.
    \end{itemize}
    \textbf{Questions:}
    \begin{enumerate}
        \item Calculate the total labor surplus (difference between marginal productivity and wages) for the platform.
        \item Determine the optimal number of workers the platform should employ to maximize total labor surplus.
    \end{enumerate}
\end{frame}

% Solution 3 Slide
\begin{frame}{Solution 3: Productivity and Labor Surplus}
    \textbf{Given Data:}
    \begin{itemize}
        \item Marginal productivity: \( MP = 50 - 2N \).
        \item Wage = \$20/hour.
        \item Current workers: \( N = 10 \).
    \end{itemize}

\end{frame}

% Solution 3 Slide
%\begin{frame}{Solution 3: Productivity and Labor Surplus}

  %  \textbf{Solution:}
 %   \begin{align*}
  %      \text{Total productivity} &= \sum_{N=1}^{10} (50 - %2N) = 50 + 48 + 46 + \cdots + 32 = 410 \\
 %       \text{Total wages paid} &= 10 \times 20 = 200 \\
    %    \text{Labor surplus} &= 410 - 200 = \$210.
   % \end{align*}
   % \textbf{Optimal Workers:}
  %  \begin{align*}
   %     \text{Set } MP = 20: & \quad 50 - 2N = 20 \\
  %      2N &= 30 \quad \Rightarrow \quad N = 15.
  %  \end{align*}
  %  \text{Optimal number of workers: \( N = 15 \).}
%\end{frame}


\begin{frame}
\frametitle{Solution 3: Productivity and Labor Surplus}

Labor surplus is the difference between the marginal productivity and the wage paid to the workers. For each worker, we calculate the difference between the marginal productivity and the wage, then sum it for all workers employed.

\textbf{Step 1: Calculate marginal productivity for \(N = 10\):}

\[
MP(10) = 50 - 2(10) = 50 - 20 = 30
\]

\textbf{Step 2: Calculate labor surplus per worker:}

\[
\text{Labor Surplus per Worker} = MP - W = 30 - 20 = 10
\]


\end{frame}



\begin{frame}
\frametitle{Solution 3: Productivity and Labor Surplus}


\textbf{Step 3: Calculate total labor surplus for 10 workers:}

\[
\text{Total Labor Surplus} = 10 \times 10 = 100
\]

Thus, the total labor surplus for the platform when employing 10 workers is \(\boxed{100}\).
\end{frame}





\begin{frame}
\frametitle{Solution 3: Productivity and Labor Surplus}

To maximize the total labor surplus, we need to maximize the total difference between the marginal productivity of labor and the wage. The total labor surplus is:

\[
\text{Total Labor Surplus}(N) = N \times (MP(N) - W)
\]
Substitute the expression for \(MP(N)\) and \(W\):

\[
\text{Total Labor Surplus}(N) = N \times \left( (50 - 2N) - 20 \right) = N \times (30 - 2N)
\]

\textbf{Step 1: Expand the equation:}

\[
\text{Total Labor Surplus}(N) = 30N - 2N^2
\]


\end{frame}



\begin{frame}
\frametitle{Solution 3: Productivity and Labor Surplus}



\textbf{Step 2: Maximize the total labor surplus by differentiating:}

\[
\frac{d}{dN}(30N - 2N^2) = 30 - 4N
\]

Set the derivative equal to 0:

\[
30 - 4N = 0
\]
\[
4N = 30
\]
\[
N = 7.5
\]

Since \(N\) must be an integer, we round to the nearest whole number, so the optimal number of workers is \(\boxed{8}\).
\end{frame}





% Problem 4 Slide
\begin{frame}{Problem 4: Gender Pay Gap in Gig Platforms}
    A study of a ride-hailing platform shows the following data:
    \begin{itemize}
        \item Male drivers work an average of 40 hours per week and earn \$1,200/week.
        \item Female drivers work an average of 30 hours per week and earn \$900/week.
        \item Controlling for hours worked, male drivers complete 5 rides per hour, while female drivers complete 4 rides per hour on average.
        \item Each ride generates \$10 in revenue.
    \end{itemize}
    \textbf{Questions:}
    \begin{enumerate}
        \item Calculate the earnings gap between male and female drivers on an hourly basis.
        \item Determine how much of the gap is due to productivity differences (rides per hour) versus potential discrimination.
    \end{enumerate}
\end{frame}


% Solution Slide for Problem 5
\begin{frame}{Solution: Gender Pay Gap in Gig Platforms}
    \textbf{1. Earnings Gap on an Hourly Basis:}
    \begin{itemize}
        \item Male drivers: 
        \[
        \text{Hourly earnings} = \frac{\$1200}{40} = \$30 \, \text{per hour}
        \]
        \item Female drivers: 
        \[
        \text{Hourly earnings} = \frac{\$900}{30} = \$30 \, \text{per hour}
        \]
        \item \textbf{Earnings Gap:} On an hourly basis, there is no earnings gap (\$30/hour for both).
    \end{itemize}
    

\end{frame}



% Solution Slide for Problem 5
\begin{frame}{Solution: Gender Pay Gap in Gig Platforms}

    
    \textbf{2. Analyzing the Gap: Productivity vs. Discrimination}
    \begin{itemize}
        \item Productivity difference (rides/hour):
        \[
        \text{Male drivers} = 5 \, \text{rides/hour}, \quad \text{Female drivers} = 4 \, \text{rides/hour}
        \]
        \item Revenue per ride = \$10.
        \item \textbf{Male's hourly revenue:} 
        \[
        5 \, \text{rides/hour} \times 10 = \$50 \, \text{per hour}
        \]
        \item \textbf{Female's hourly revenue:}
        \[
        4 \, \text{rides/hour} \times 10 = \$40 \, \text{per hour}
        \]
        \item \textbf{Conclusion:} The productivity difference contributes \$10 per hour to the earnings gap, but no gap remains after controlling for hours worked.
    \end{itemize}
\end{frame}




\end{document}

